% Requires amsmath, amssymb, booktabs, graphicx and url (or hyperref).
% Copy both figure PDFs to Images/ and retain FROST_hardware_tools.bib entries.
\section{Hardware resource savings through categorical value selection}
\label{sec:fpga_evaluation}

FROST reduces the arithmetic resources of attention aggregation by
replacing dense probability--value products with categorical value
selection. In the implemented fused kernel, this mapping required no
dedicated digital signal processing (DSP) blocks or block random-access
memory (BRAM) instances. Its field-programmable gate array (FPGA)
combinational resource count was 71.3\% lower than the four-lane
single-precision floating-point (FP32) control and 26.5\% lower than the
single-lane FP32 control. These savings were obtained with random-number
generation, sampling-state updates, coefficient storage, buffering and
format conversion included in the implementation boundary
(Fig.~\ref{fig:fused_hardware_stage3}).
In a separate standard-cell assessment, the short-stream operating point
also reduced mapped area, estimated active power and tile latency
simultaneously relative to four-lane FP32 at the same clock frequency.

\paragraph{Architectural basis of the resource savings.}
The kernel maps 16 attention scores and a value matrix
$\boldsymbol V\in\mathbb{R}^{16\times16}$ to a 16-component output.
For stream length $L$, FROST forms the sampled estimate
\begin{equation}
\widehat{\boldsymbol y}
=\frac{1}{L}\sum_{t=1}^{L}\boldsymbol V_{J_t,:},
\label{eq:fused_categorical_attention}
\end{equation}
where $t$ indexes retained samples, $J_t\in\{1,\ldots,16\}$ is the
category emitted by the tensorized finite-state operator, and
$\widehat{\boldsymbol y}\in\mathbb{R}^{16}$ is the estimated attention
output. Each sample selects one value row, which is added to sixteen
signed 32-bit accumulators. This formulation removes explicit probability
materialization and dense probability--value multiplication from the
aggregation datapath. For the supported power-of-two stream lengths,
final scaling is implemented by binary shifting.

The frozen hardware candidate uses tensor-train rank 2, 16
finite-state-machine (FSM) states per score input, 8-bit cumulative
distribution function (CDF) entries and 16 burn-in updates.
Fifteen binary latent transitions each retain one non-terminal threshold;
the final category transition retains 15 thresholds. Making terminal
unit-CDF entries implicit leaves 944 bytes of constant lookup data before
logic mapping. Sorting, inverse index mapping and 32 separately
initialized xorshift32 states are included. A single circuit supports
$L\in\{64,256,1024\}$: the stream length changes execution duration
without changing the allocated resources.

\paragraph{Matched interfaces and implementation controls.}
We implemented FP32 controls with one and four shared fused
multiply-add (FMA) lanes using Berkeley
HardFloat~\cite{hauser2019hardfloat}, together with a four-lane
fixed-point control (FX16). Both FP32 controls use a 1,025-entry
exponential look-up table (LUT), spaced at $1/64$ over score differences
from 0 to 16. They accumulate unnormalized exponential--value products
and use one reciprocal per query to normalize the final vector.
FX16 uses a LUT with 16 fractional bits, widened accumulators and one
restoring reciprocal per query.

All implementations receive signed 16-bit scores and values with ten
fractional bits and return signed 32-bit outputs with ten fractional
bits. FP32 computation uses the integer-unit value representation before
conversion back to this common interface. A tile contains 16 consecutive
queries sharing one $16\times16$ value matrix. Value loading therefore
occurs once per tile; each query loads its own 16 scores. Transfers use
one word per cycle, and queries do not overlap. Burn-in and output
transfers are included; one-time random-seed initialization is excluded.
The comparison isolates hardware cost without matching approximation
error. Query--key products, projections, feed-forward layers and
whole-network classification are outside the evaluated kernel.

\paragraph{FPGA resource allocation.}
We targeted an LFE5U-85F FPGA with CABGA381 packaging and speed grade 6,
using Yosys 0.33 and nextpnr-ecp5
0.6~\cite{yosysSoftware,nextpnrSoftware}. All four implementations were
routed with seeds 1, 2 and 3 at a common 8-MHz clock, with identical
resource counts across seeds. This common frequency was selected after
the initial FP32 implementation failed the 15-MHz constraint; all four
designs were then rerouted at 8 MHz. Board-specific pin locations and
external input/output delays were not constrained.

Table~\ref{tab:fused_fpga_resources} shows the complete resource vector.
Relative to the four-lane FP32 control, FROST removes all 16 DSP and
three BRAM instances and reduces combinational resources from 23,799
to 6,834. The same DSP- and BRAM-free mapping also distinguishes it
from FX16, which uses seven DSP and two BRAM instances. FROST instead
uses distributed value storage and constant logic-mapped CDFs; zero
BRAM usage does not imply zero storage cost. Its 2,699 flip-flops (FFs)
exceed the 2,043 of either FP32 control, while FX16 uses fewer
combinational resources. The demonstrated gain is therefore reduced
dedicated-block demand and lower combinational cost relative to FP32,
with the state-storage cost explicitly retained.

\begin{table}[!htbp]
\centering
\caption{Post-route resource allocation for the fused attention kernel.}
\label{tab:fused_fpga_resources}
\small
\setlength{\tabcolsep}{4pt}
\begin{tabular*}{\linewidth}{@{\extracolsep{\fill}}lrrrr@{}}
\toprule
Implementation & COMB & FF & DSP & BRAM \\
\midrule
FP32-LUT, 1 lane & 9,304 & 2,043 & 4 & 3 \\
FP32-LUT, 4 lanes & 23,799 & 2,043 & 16 & 3 \\
FX16, 4 lanes & 3,419 & 1,280 & 7 & 2 \\
FROST, all supported $L$ & 6,834 & 2,699 & 0 & 0 \\
\bottomrule
\end{tabular*}
\par\smallskip\footnotesize
COMB denotes ECP5 \texttt{TRELLIS\_COMB}, including arithmetic and
distributed-memory logic, rather than a literal LUT4 count. DSP and
BRAM denote \texttt{MULT18X18D} and nominal 18-kbit \texttt{DP16KD}
instances. The FROST circuit supports $L=64$, 256 and 1024.
All counts include the declared kernel interfaces and auxiliary logic.
\end{table}

\begin{figure}[!htbp]
\centering
\includegraphics[width=\linewidth]{Images/FROST_hardware_stage3.pdf}
\caption{\textbf{Hardware resource substitution through categorical value selection.}
\textbf{a}, Functional organization of the fused kernel. Conventional
controls accumulate products of exponential coefficients and values;
FROST selects value rows using the tensorized FSM sampler and accumulates
the selected values. Shared input and output formats are retained.
The schematic groups operations rather than depicting physical placement.
\textbf{b}, Absolute post-route resource counts on the same ECP5 device.
Axes are independent across resource classes, and all begin at zero.
FROST uses one programmable circuit for all three stream lengths.
Its zero DSP and BRAM counts coexist with distributed storage, constant
CDF logic and increased FF usage relative to FP32.}
\label{fig:fused_hardware_stage3}
\end{figure}

\paragraph{Resource savings at a short-stream operating point.}
At $L=64$, FROST completes a 16-query tile in 2,592 cycles, compared
with 2,816 cycles for four-lane FP32. At the common 8-MHz FPGA clock,
these correspond to 324 and 352~$\mu$s, respectively. Thus, the
71.3\% reduction in combinational resources is accompanied by an
8.0\% reduction in tile latency, rather than an increase in execution
time at this operating point. Increasing $L$ retains the same resource
allocation while increasing the tile cycle count
(Table~\ref{tab:fused_cycles}).

\begin{table}[!htbp]
\centering
\caption{Tile execution time across the supported stream lengths.}
\label{tab:fused_cycles}
\small
\setlength{\tabcolsep}{4pt}
\begin{tabular*}{\linewidth}{@{\extracolsep{\fill}}lrr@{}}
\toprule
Implementation & Tile cycles & Tile latency at 8 MHz ($\mu$s) \\
\midrule
FP32-LUT, 1 lane & 6,080 & 760 \\
FP32-LUT, 4 lanes & 2,816 & 352 \\
FX16, 4 lanes & 2,848 & 356 \\
FROST, $L=64$ & 2,592 & 324 \\
FROST, $L=256$ & 5,664 & 708 \\
FROST, $L=1024$ & 17,952 & 2,244 \\
\bottomrule
\end{tabular*}
\par\smallskip\footnotesize
Each tile includes loading 256 value words, loading 16 scores for each
of 16 queries, computation, burn-in where applicable, and all output
transfers. Cycles were obtained from register-transfer-level simulation.
These are attention-kernel tile latencies, not image-inference latencies.
\end{table}

\paragraph{Independent standard-cell assessment.}
The same register-transfer-level (RTL) designs were mapped to the
Nangate45 typical standard-cell library at 1.10 V and $25\,^{\circ}$C.
All memories were mapped into standard cells. OpenSTA 3.1.0 estimated
power from zero-delay gate-level switching activity at a common
15-MHz clock~\cite{openstaSoftware}, using the first 16-query evaluation
tile and three stream seeds for FROST. All reported designs met
pre-layout register timing at this frequency.

Let $C_{\mathrm{tile}}$ be the tile cycle count, $f$ the clock frequency
in hertz, and $P_{\mathrm{cell}}$ the estimated active cell power in watts,
including internal, switching and leakage components. Estimated tile
energy $E_{\mathrm{tile}}$, in joules, is
\begin{equation}
E_{\mathrm{tile}}=P_{\mathrm{cell}}\frac{C_{\mathrm{tile}}}{f}.
\label{eq:fused_tile_energy}
\end{equation}
The standard-cell results use $f=15$ MHz throughout; their power is
not combined with the FPGA timing results.

At $L=64$, FROST's mapped area is
$63.62\times10^3~\mu\mathrm{m}^2$, compared with
$85.36\times10^3~\mu\mathrm{m}^2$ for four-lane FP32, a 25.5\%
reduction. Estimated active power decreases from 3.556 to 2.650~mW,
a 25.5\% reduction. At the common 15-MHz clock, the corresponding
tile latencies are 187.7 and 172.8~$\mu$s, an 8.0\% reduction.
FROST therefore improves all three power, performance and area (PPA)
metrics relative to this four-lane FP32 implementation at $L=64$
within the same standard-cell flow. Estimated energy per tile decreases
from 0.668 to 0.458~$\mu$J, a 31.4\% reduction
(Fig.~\ref{fig:fused_hardware_tradeoffs}).
At $L=256$, estimated energy remains 16.0\% below single-lane FP32
but exceeds four-lane FP32 and FX16, showing the execution cost of
longer sampling at fixed resource allocation.

\begin{table}[!htbp]
\centering
\caption{Standard-cell area, power and latency at a common 15-MHz clock.}
\label{tab:fused_asic_power}
\small
\setlength{\tabcolsep}{3pt}
\begin{tabular*}{\linewidth}{@{\extracolsep{\fill}}lrrrr@{}}
\toprule
Implementation
& \shortstack{Area\\($10^3\,\mu\mathrm{m}^2$)}
& \shortstack{Active power\\(mW)}
& \shortstack{Tile latency\\($\mu$s)}
& \shortstack{Energy/tile\\($\mu$J)} \\
\midrule
FP32-LUT, 1 lane & 62.86 & 3.598 & 405.3 & 1.458 \\
FP32-LUT, 4 lanes & 85.36 & 3.556 & 187.7 & 0.668 \\
FX16, 4 lanes & 60.08 & 2.839 & 189.9 & 0.539 \\
FROST, $L=64$ & 63.62 & $2.650\pm0.004$ & 172.8 & $0.458\pm0.001$ \\
FROST, $L=256$ & 63.62 & $3.244\pm0.010$ & 377.6 & $1.225\pm0.004$ \\
FROST, $L=1024$ & 63.62 & $3.619\pm0.012$ & 1196.8 & $4.331\pm0.015$ \\
\bottomrule
\end{tabular*}
\par\smallskip\footnotesize
All entries are pre-layout tool estimates for Nangate45 at 15 MHz,
not FPGA power or board measurements. Area is summed mapped-cell area.
FROST values are means $\pm$ sample standard deviations over three
stream seeds on one activity tile; deterministic controls use that tile
once. Clock-tree implementation, extracted interconnect, time-resolved
glitches, pads and external memory are excluded.
\end{table}

\begin{figure}[!htbp]
\centering
\includegraphics[width=\linewidth]{Images/FROST_hardware_tradeoffs.pdf}
\caption{\textbf{Simultaneous PPA reductions at the short-stream operating point.}
Both panels use the same Nangate45 standard-cell flow at 15 MHz.
\textbf{a}, FROST at $L=64$, normalized separately by the mapped area,
estimated active power, tile latency and estimated tile energy of
four-lane FP32. Lower values are preferable; the dashed line denotes
the FP32 reference. Percentage reductions are calculated from unrounded
results. \textbf{b}, Estimated FROST energy across all supported stream
lengths, with deterministic controls shown as horizontal references.
Error bars show sample standard deviations over three stream seeds on
one tile and can be smaller than the symbols. Each tile contains 16
queries. Values are pre-layout estimates without matched approximation
error.}
\label{fig:fused_hardware_tradeoffs}
\end{figure}

\paragraph{Verification and interpretation.}
Verilator RTL simulations covered 3,200 query rows from the 200-image
evaluation subset. The three FROST stream lengths and stream seeds
4004, 5005 and 6006, together with the deterministic controls, produced
614,400 output components that agreed bit for bit with independent
integer or C++ floating-point references. Gate-level simulations matched
the RTL for all 3,072 checked components, and every power analysis
reported zero unannotated pins.
These checks establish implementation agreement for the evaluated inputs.
The resource result is a realizable DSP- and BRAM-free mapping of the
fused sampling kernel, with lower combinational cost than both FP32
controls. At $L=64$, the standard-cell implementation also reduces area,
estimated power and latency simultaneously relative to four-lane FP32.
These findings concern
the declared kernel and do not establish whole-Transformer or
quality-matched hardware superiority.
